\documentclass[11pt,letterpaper]{article}
\pdfinfoomitdate=1
\pdftrailerid{}
\pdfmapfile{+lm.map}
\pdfmapfile{+cm.map}
\pdfmapfile{+symbols.map}
\usepackage[T1]{fontenc}
\usepackage[utf8]{inputenc}
\usepackage{lmodern,amsmath,amssymb,amsthm,mathtools}
\usepackage[margin=1in]{geometry}
\usepackage{microtype,booktabs,longtable,array,enumitem}
\usepackage{xurl}
\usepackage[hidelinks,pdfauthor={},pdftitle={Moving fugacity asymptotics for distinct partitions}]{hyperref}
\usepackage{fancyhdr,needspace}
\pagestyle{fancy}\fancyhf{}\fancyhead[L]{\small Moving fugacity for distinct partitions}\fancyhead[R]{\small A291698 and A292304}\fancyfoot[C]{\thepage}
\renewcommand{\headrulewidth}{0pt}
\setlength{\headheight}{14pt}
\setlength{\parindent}{0pt}\setlength{\parskip}{5pt}
\setlist{nosep,leftmargin=1.5em}
\providecommand{\tightlist}{\setlength{\itemsep}{0pt}\setlength{\parskip}{0pt}}
\allowdisplaybreaks[2]
\setcounter{tocdepth}{2}
% [write] notes added when the report entered the collection (batch 77)
\newenvironment{writenote}[1][]{\par\smallskip\begingroup\small
 \noindent\textbf{[write]}\if\relax\detokenize{#1}\relax\else~\textbf{#1.}\fi\ }%
 {\par\endgroup\smallskip}
\title{\textbf{Moving fugacity asymptotics for distinct partitions}\\[7pt]\large Uniform coefficient expansions and inverses for A291698 and A292304}
\author{}\date{1 October 2026}
\begin{document}
\maketitle\thispagestyle{empty}
\begin{abstract}
For distinct partitions with size-dependent fugacity $u=n^\alpha$, this report proves a uniform all-orders Bessel expansion on every compact range $0<a\le\alpha\le b$. An exact Jacobi transformation identifies finitely many conjugate exponential sectors, with rigorously separated local sector integrals and resonance-safe remainder bounds. The results imply the displayed leading equivalents for OEIS A291698 ($\alpha=1$) and A292304 ($\alpha=2$). A specified real continuation supports a controlled inverse, beginning with Lambert $W$ and an explicit logarithmic correction. Exact computations expose a substantial preasymptotic oscillation: at $n=20{,}000$, the principal approximation for $\alpha=2$ still differs by about $5.6\%$. The scope is fixed sector count and fixed algebraic depth; no infinite Bessel-sector sum, growing-sector theorem, or numerical certification from unspecified asymptotic constants is claimed.
\end{abstract}
\vspace{5pt}
\textbf{Reading guide.} Sections 1--5 contain the coefficient and sector proofs. Section 6 gives the two OEIS specializations, Section 7 the inverse, and Sections 8--10 the computations, literature boundary, and further questions. The appendix records reproducibility details. The report makes no claim of publication priority or external peer review.
\newpage
\tableofcontents
\newpage
\section{Result and scope}\label{mfp:result-and-scope}

Let

\[
 F(q,u)=\prod_{k\ge1}(1+u q^k),\qquad
 a_n(\alpha)=[q^n]F(q,n^\alpha),\qquad 0<a\le\alpha\le b<\infty.
\]

The endpoints \(a,b\) are fixed. Put

\[
 \begin{gathered}
 u=n^\alpha,\quad L=\log u,\quad A=-\operatorname{Li}_2(-u),\quad c_1(u)=\frac{u}{12(1+u)},\\
 N=n+c_1(u),\qquad t=\sqrt{A/N},\qquad s=\sqrt{AN}.
 \end{gathered}
\]

Thus \(t\) is comparable to \(\log n/\sqrt n\), \(s\) to
\(\sqrt n\log n\), uniformly in \(\alpha\). All square roots below are
principal; for fixed integer \(j\) they have positive real part once
\(n\) is large.

Let \(D=u\,d/du\) and define

\[
 c_r(u)=\frac{B_{2r}}{(2r)!}D^{2r-1}\log(1+u),\qquad
 \sum_{m\ge0}h_m(u)z^m
 =\exp\left(\sum_{r\ge2}c_r(u)z^{2r-1}\right)
\]

as a \textbf{formal} power series. No convergence of the infinite
Euler--Maclaurin series is asserted. Define

\[
 A_j=A-2\pi^2j^2+2\pi i jL,\qquad
 T_{j,m}=\left(\frac{A_j}{N}\right)^{(m+1)/2}
 I_{m+1}(2\sqrt{NA_j}),
\]

and

\[
 B_0=(1+u)^{-1/2}tI_1(2s),\qquad
 \Delta_j=2s-2\Re\sqrt{NA_j}.
\]

\subsection{Theorem 1 Principal sector expansion}\label{mfp:principal-sector-theorem}

For every fixed integer \(M\ge0\),

\[
 a_n(\alpha)=(1+u)^{-1/2}\sum_{m=0}^{M}h_m(u)t^{m+1}I_{m+1}(2s)
 +B_0 O_{a,b,M}\left(\frac{t^{M+1}}{u}
 +\exp\{-c/[t(1+L)^2]\}\right). \tag{1}
\]

Here \(c>0\) depends only on \(a,b\) (indeed the underlying interval
bound is absolute). The second error is smaller than every power of
\(n\), uniformly in \(\alpha\). Consequently it can be absorbed into the
first error asymptotically. In particular, taking \(M=2\) gives

\[
 a_n(\alpha)=B_0\{1+O_{a,b}(t^3/u)\}. \tag{2}
\]

The absorbed form of (2) can have a very late numerical onset. The
explicit exponential remainder in (1) should be retained in practical
work.

\subsection{Theorem 2 Finite resonance sectors}\label{mfp:finite-resonance-sector-theorem}

For fixed nonnegative integers \(K,M\), the stronger, finite-sector
statement is

\[
 a_n(\alpha)=(1+u)^{-1/2}
 \sum_{|j|\le K}(-1)^j\sum_{m=0}^{M}h_m(u)T_{j,m}
 +B_0 O_{a,b,K,M}\left(\frac{t^{M+1}}{u}+e^{-\Delta_{K+1}}\right). \tag{3}
\]

The \(j\) and -\(j\) terms are conjugates, so the displayed approximant
is real. For each fixed \(j\ge1\),

\[
 \Delta_j\sim\frac{2\sqrt2\pi^4}{3}\frac{j^2\sqrt n}{L^3}
 \sim\frac{2\pi^4j^2}{3tL^2}. \tag{4}
\]

Uniformity in (3) is for \textbf{fixed \(K,M\)} and \(\alpha\) in a
fixed positive compact interval. Neither \(K\) growing with \(n\) nor an
infinite sum of Bessel sectors is asserted. Section 5 gives the
finite-arc proof, including the choice of an auxiliary larger fixed
sector cutoff. The exact local-sector version is stated separately in
Section 5.4.

Formula (3) is an all-algebraic-order expansion inside finitely many
exponentially small sectors. At any fixed \(M\), its principal-sector
algebraic remainder ultimately exceeds all retained exponential sectors;
it is not a claim of exponentially precise hyperasymptotics with \(M\)
growing with \(n\). The exact local-sector formulation in Section 5.4
supplies a stronger, non-vacuous exponential separation before the
algebraic expansions are truncated.

\subsection{Provenance, status and reading conventions}\label{mfp:provenance}

\begin{writenote}[Added 2 October 2026, batch~77]
\emph{Provenance.} This report was built from one manuscript: batch~77,
manuscript~34 of the repository's incoming-reports channel, archive
\nolinkurl{moving-fugacity-reproducibility.zip} (main file
\nolinkurl{moving-fugacity-report.tex}, 985 lines, dated 1 October 2026,
with a 17-page PDF that is not shipped), arrival commit
\texttt{096ee7b87}, placed in \texttt{4f11bc9c0}. The text from the title
to the end of the scope checklist is the delivered manuscript, unchanged
except for the label prefix \texttt{mfp:} and the notes marked
\textbf{[write]}; no statement, proof, number or symbol of it was changed,
and every section number the text cites by hand is still correct. The delivered proof note \nolinkurl{theorem-and-proof.md}
and the delivered \nolinkurl{REPRODUCIBILITY.md} are shipped beside this
article, byte-identical; the SHA-256 printed in Appendix~A still matches
the shipped \nolinkurl{theorem-and-proof.md}. \emph{Pin.} The delivery
names no ProveIt revision and continues no repository report.
\emph{Status.} An unrefereed research manuscript (its author line is
empty). Nothing in it is formalized in Lean or Rocq, and its placement in
the repository confers no formal status. At intake and at this write all
four delivered programs (Appendix~A) were rerun on a copy, with outputs
identical to the shipped records up to line endings; the first ten terms
of both OEIS sequences were recomputed independently by a direct product
expansion. These are numerical diagnostics; the proofs were not
re-derived.
\end{writenote}

\begin{writenote}[Reading conventions]
The manuscript reuses several letters; no symbol was renamed. Two readings
need care. \emph{The fugacity $u=n^\alpha$ is one weight per part, the
same for every part of a partition of $n$}: a partition with $m$ parts has
weight $n^{\alpha m}$. It is \emph{not} a part-dependent weight $k^\alpha$
on a part of size $k$, which is the different family $\prod_k(1+k^\alpha
q^k)$ (see the note at the end of Section~9). \emph{The symbol $C_n$ has
two unrelated meanings}: in Section~2 it is the Bessel scale
$C_n=A^{1/4}e^{2\sqrt{AN}}/(2\sqrt\pi N^{3/4}\sqrt{1+u})$, while in
Sections~5.2--5.3 it is the bounded constant $C_n=A-L^2/2\to\pi^2/6$.

\smallskip
\noindent\begin{tabular}{@{}>{\raggedright\arraybackslash}p{0.17\linewidth}
 >{\raggedright\arraybackslash}p{0.77\linewidth}@{}}
\toprule
symbol & meanings, by section \\
\midrule
$a,b$ & the endpoints $0<a\le\alpha\le b$ of the compact $\alpha$-range;
 $a_n(\alpha)$ is the sequence; $a_j,b_j$ are $\Re\sqrt{A_j}$,
 $\Im\sqrt{A_j}$ in Section~5.3; $b_r(\nu)$ are Bessel coefficients
 (Section~2). \\
$A$ & $A=-\operatorname{Li}_2(-u)$; $A_j$ its sector shifts; $A_\infty$ its
 elementary part (Section~6); $\mathcal A_\alpha(x)$ the real continuation
 (13). \\
$B$ & $B_{2r}$ Bernoulli numbers; $B_0$ the principal scale; $\mathcal B_j$
 the exact sector integrals (12c). \\
$c$ & $c_r(u)$ endpoint coefficients ($c_1(u)=u/(12(1+u))$); $c$, $c_0$,
 $c_J$, $c_j$ unnamed positive constants, and $c_1$ in (9) is such a
 constant, not $c_1(u)$. \\
$v$ & $v=1/u$, the small fugacity (Sections~5.1, 7); a summation index in
 $b_r(\nu)$ and in the residue series after (10). \\
$y$ & $y=tx-L$ in Section~3; $y=\theta/t$ in Section~5.2; $y=\log X$ in
 Section~7. \\
$z$ & the complex variable $q=e^{-z}$; but $z=y/(2\sqrt2\alpha)$ in (15). \\
$m$, $r$ & $m$ is the endpoint degree, and also the number of consecutive
 indices in (9); $r$ indexes $c_r$, is the starting index in (9) and the
 Lagrange order in (17); $r_0$ in (15). \\
$x$ & the real continuation variable of Section~7; an integration variable
 ($w=xz$) in Section~3. \\
$D$, $E$, $g$, $H$ & $D=u\,d/du$ (and $D_v$); $E(x)=\log S_{K,M}(x)-g(x)$
 in (17), while $E_n(\alpha)$ is the Sommerfeld expression of Section~6;
 $g(x)=\sqrt2\alpha\sqrt x\log x$, while $g_d$ are generator coefficients
 (Section~2); $H=g^{-1}$, while $H_{\rm ex}$ is the amplitude (12b). \\
$\eta$ & a small fixed constant in (5); $\eta(y)=E(H(y))$ in (17); the
 Dedekind-$\eta$ transformation in Section~5. \\
$R_K$, $R_P$ & the exact remainder of (12d); the Euler--Maclaurin remainder
 of (13a). \\
$t$, $s$, $N$, $L$ & $t=\sqrt{A/N}$ (the saddle, not a time), $s=\sqrt{AN}$,
 $N=n+c_1(u)$, $L=\log u$. \\
\bottomrule
\end{tabular}
\end{writenote}

\section{Coefficient generator and ordinary relative
expansion}\label{mfp:coefficient-generator-and-ordinary-relative-expansion}

The first endpoint coefficients are

\[
 c_2=\frac{u(u-1)}{720(1+u)^3},\qquad
 c_3=-\frac{u(u-1)(u^2-10u+1)}{30240(1+u)^5},
\]

\[
 c_4=\frac{u(u-1)(u^4-56u^3+246u^2-56u+1)}{1209600(1+u)^7}.
\]

For each fixed \(r\ge2\), \(c_r(u)=O_r(u^{-1})\). The \(h_m\) generator
can be evaluated without symbolic exponentiation. Set \(g_{2r-1}=c_r\)
for \(r\ge2\) and all other \(g_d=0\); then

\[
 h_0=1,\qquad m h_m=\sum_{d=1}^m d g_d h_{m-d}.
\]

In particular,

\[
 h_1=h_2=h_4=0,\quad h_3=c_2,\quad h_5=c_3,\quad
 h_6=c_2^2/2,\quad h_7=c_4,\quad h_8=c_2c_3,\quad
 h_9=c_5+c_2^3/6.
\]

For every fixed \(m>0\), \(h_m=O_m(u^{-1})\). This last bound alone does
not prove a remainder with \(1/u\); that issue is addressed explicitly
in Section 3.

The ordinary large-argument Bessel expansion is

\[
 I_\nu(2s)\sim\frac{e^{2s}}{\sqrt{4\pi s}}
 \sum_{r\ge0}b_r(\nu)s^{-r},\qquad
 b_r(\nu)=\frac{(-1)^r}{r!16^r}\prod_{v=1}^r[4\nu^2-(2v-1)^2].
\]

Thus the principal relative expansion is generated by the double series

\[
 \frac{a_n}{C_n}\sim\sum_{m,r\ge0}h_m(u)t^m b_r(m+1)s^{-r},\qquad
 C_n=\frac{A^{1/4}e^{2\sqrt{AN}}}{2\sqrt\pi\,N^{3/4}\sqrt{1+u}}.
\]

Its first Bessel terms are
\[1-\frac{3}{16s}-\frac{15}{512s^2}-\frac{105}{8192s^3}.\] The first new
endpoint contribution is \(c_2t^3I_4(2s)/I_1(2s)\), of order
\(n^{-\alpha-3/2}\log^3n\). Keeping exact Bessels is preferable to
expanding them numerically. For any desired fixed algebraic precision
\(n^{-R}\), choose \(M\) and the Bessel expansion depth sufficiently
large; the log factors are accounted for by choosing strict inequalities
in the depth.

\section{Euler--Maclaurin and its uniform
remainder}\label{mfp:eulermaclaurin-and-its-uniform-remainder}

Write \(f_u(w)=\log(1+ue^{-w})\), with the branch continued from the
positive real axis. Fix a small \(\eta>0\). In the strip-like sector

\[
 \Re z=t>0,\qquad |\Im z|\le\eta t/(1+L), \tag{5}
\]

the path \(w=xz,\ x\ge0\), stays uniformly away from the zeros of
\(1+ue^{-w}\). Near \(tx=L\) its argument is bounded by \(\eta\); away
from that transition its modulus is either much larger or much smaller
than one. For each integer \(r\ge2\),

\[
 \int_0^\infty |f_u^{(r)}(xz)|\,dx\le C_r/t.
\]

One direct verification uses \(y=tx-L\): the derivatives are rational
functions of \(ue^{-xz}\), bounded by \(C_re^{-|y|}\) outside a fixed
\(y\) interval and bounded uniformly inside it. The integral term is
\(A/z\), by the analytic antiderivative or contour rotation.
Euler--Maclaurin through index \(P\) therefore gives

\[
 \log F(e^{-z},u)=A/z-\tfrac12\log(1+u)
 +\sum_{r=1}^{P}c_r(u)z^{2r-1}+O_P(t^{2P-1}), \tag{6}
\]

uniformly in (5). The generic remainder here does \textbf{not} contain
\(1/u\).

To infer (1), fix \(M\) and take \(P\) with \(2P>M+2+2b\). Then

\[
 t^{2P-1}=O_{a,b,M,P}(t^{M+1}/u),
\]

because \(u\le n^b\) and \(t\) is comparable to \(\log n/\sqrt n\). Each
discarded endpoint term with degree at least \(M\)+1 is individually
\(O(t^{M+1}/u)\), since \(c_r=O(u^{-1})\). Exponentiating and collecting
coefficients through degree \(M\) gives

\[
 F(e^{-z},u)=\frac{e^{A/z+c_1z}}{\sqrt{1+u}}
 \left\{\sum_{m=0}^{M}h_m(u)z^m+O(t^{M+1}/u)\right\}. \tag{7}
\]

For \(M=0,1,2\) the first nonzero omitted endpoint degree is 3, yielding
the sharper \(O(t^3/u)\) if desired. An alternative derivation using the
exact identity in Section 5 obtains the \(1/u\) factor directly from the
small-fugacity factor's derivative norm.

\section{Minor arcs and Bessel
transfer}\label{mfp:minor-arcs-and-bessel-transfer}

The Cauchy integral on \(q=e^{-t+i\theta}\) is

\[
 a_n=\frac1{2\pi}\int_{-\pi}^{\pi}e^{n(t-i\theta)}F(e^{-t+i\theta},u)\,d\theta.
\]

Set \(p_k=ue^{-kt}/(1+ue^{-kt})\). Exactly,

\[
 \left|\frac{F(e^{-t+i\theta},u)}{F(e^{-t},u)}\right|
 \le\exp\left\{-2\sum_{k\ge1}p_k(1-p_k)\sin^2(k\theta/2)\right\}. \tag{8}
\]

On the consecutive integer interval \(|kt-L|\le1\) there are m
comparable to \(1/t\) indices, all with \(p_k(1-p_k)\ge c_0>0\). The
following elementary bound, for \(m\ge2\), is independent of the
starting index:

\[
 \sum_{k=r}^{r+m-1}\sin^2(k\theta/2)
 \ge c_1m\min\{m^2\theta^2,1\},\qquad |\theta|\le\pi. \tag{9}
\]

Indeed the minimum over the initial phase is at least

\[
 m/2-\tfrac12\left|\frac{\sin(m\theta/2)}{\sin(\theta/2)}\right|.
\]

For \(m|\theta|\le1\), Taylor bounds on sin give a fixed multiple of
\(m^3\theta^2\). For \(1\le m|\theta|\le C\), compactness of
\(|\sin x|/|x|\) away from zero gives a fixed fraction of m, uniformly
as m grows. For \(m|\theta|\ge C\), choose \(C>2\pi\) and use
\(|\sin(\theta/2)|\ge|\theta|/\pi\). This proves (9), with finitely many
small \(m\ge2\) absorbed by reducing the constant. The exclusion \(m=1\)
is necessary, and irrelevant here because m is comparable to \(1/t\) and
tends to infinity.

For \(|\theta|\ge\eta t/(1+L)\), (8)-(9) give a loss at least
\(c/[t(1+L)^2]\). Polynomial factors from normalization can be absorbed
by reducing c.~Since

\[
 1/[t(1+L)^2]\asymp\sqrt n/(\log n)^3,
\]

these arcs are uniformly beyond all algebraic orders. This argument
explicitly controls the near-resonances \(\theta\approx2\pi/\sqrt{2n}\);
a fixed-fugacity variance bound alone is not a substitute.

On the remaining arc, insert (7). For \(z=t+i\theta\),

\[
 \Re(A/z+Nz)=2s-\frac{s(\theta/t)^2}{1+(\theta/t)^2}.
\]

Consequently integrating a uniform \(O(t^{M+1}/u)\) error costs only
\(O(e^{2s}t/\sqrt s)\text{ times that error}\), the scale of the main
Bessel integral.

For each fixed m, deform the truncated vertical segment to the nearby
arc of \(|z|=t\), adding short radial connectors. At an endpoint
\(\theta=\eta t/(1+L)\), all connectors lose at least \(cs/(1+L)^2\),
comparable to \(c/t\), in real exponent. Completing the circle loses the
same amount. The exact residue is

\[
 \frac1{2\pi i}\oint_{|z|=t}z^m e^{A/z+Nz}\,dz
 =(A/N)^{(m+1)/2}I_{m+1}(2\sqrt{AN}). \tag{10}
\]

For completeness, expand both exponentials: the residue is
\(\sum_{v\ge0}A^{v+m+1}N^v/((v+m+1)!v!)\), which is exactly the Bessel
series. Equations (7)-(10) prove Theorem 1.

\section{Exact Jacobi decomposition and the finite-sector
proof}\label{mfp:exact-jacobi-decomposition-and-the-finite-sector-proof}

Let \(q=e^{-z},\ \Re z>0,\ u>1,\ L=\log u,\ Q=e^{-4\pi^2/z}\). With
\((a;q)_\infty=\prod_{r\ge0}(1-aq^r)\), the following is an
\textbf{exact identity}:

\[
 F(e^{-z},u)=\frac{u^{-1/2}\exp\{(L^2/2+\pi^2/6)/z+z/12\}}
 {(Q;Q)_\infty(-1/u;e^{-z})_\infty}
 \sum_{j\in\mathbb Z}(-1)^j
 \exp\{(-2\pi^2j^2+2\pi i jL)/z\}. \tag{11}
\]

Proof: Jacobi's triple product gives

\[
 F(q,u)(-1/u;q)_\infty(q;q)_\infty
 =\sum_{m\in\mathbb Z}u^m q^{m(m+1)/2}.
\]

Complete the square and apply Gaussian Poisson summation. The bilateral
sum becomes

\[
 \sqrt{2\pi/z}\,e^{L^2/(2z)-L/2+z/8}
 \sum_{j\in\mathbb Z}(-1)^j e^{(-2\pi^2j^2+2\pi ijL)/z}.
\]

The \(\eta\) transformation is

\[
 (e^{-z};e^{-z})_\infty
 =\sqrt{2\pi/z}\,e^{-\pi^2/(6z)+z/24}(Q;Q)_\infty.
\]

Division proves (11), first for \(z>0\) and then throughout
\(\Re z>0\) by analytic continuation, with consistent
square-root branches.

\subsection{Uniform local amplitude over any fixed number of
sectors}\label{mfp:uniform-local-amplitude-over-any-fixed-number-of-sectors}

Fix \(C>0\) and allow \(|\Im z|\le Ct/L,\ \Re z=t\). Unlike (6), \(C\)
need not be small. Apply Euler--Maclaurin to
\(\log(-1/u;e^{-z})_\infty\), including its \(k=0\) endpoint. Its
fugacity \(v=1/u\) tends to zero. For \(r\ge1\), the corresponding
derivative \(L^1\) norm is \(O_r(1/(ut))\), because \(|ve^{-xz}|\le v\) and
the denominator is at least \(1-v\). Hence its remainder is directly
\(O(t^{2P-1}/u)\), throughout this larger sector.

The identity

\[
 A=\tfrac12L^2+\pi^2/6+\operatorname{Li}_2(-1/u)
\]

and the relations \(1/12-c_1(1/u)=c_1(u)\), \(c_r(1/u)=-c_r(u)\) for
\(r\ge2\) transform (11) into

\[
 F(e^{-z},u)=\frac1{\sqrt{1+u}}
 \left\{\sum_{m=0}^{M}h_m(u)z^m+O(t^{M+1}/u)\right\}
 \sum_{j\in\mathbb Z}(-1)^j e^{A_j/z+c_1z}
 \{1+O(e^{-c/t})\}. \tag{12}
\]

The last factor comes from \((Q;Q)_\infty\), since \(\Re(1/z)\) is
comparable to \(1/t\) for fixed \(C\) and \(L\) tending to infinity. For
a fixed desired \(M\) choose \(2P-1\ge M+1\). The estimate on the brace
is uniform in \(z\); its error multiplies the theta sum, so
estimates below use the \textbf{sum of the absolute values} of the
theta terms, never division by a possibly vanishing theta sum.

\subsection{Isolating and completing finitely many saddle
arcs}\label{mfp:isolating-and-completing-finitely-many-saddle-arcs}

Fix an auxiliary integer \(J>K\), to be chosen sufficiently large below.
Restrict initially to

\[
 |\theta|\le(2J+1)\pi t/L.
\]

Split this interval at \(\theta=(2j+1)\pi t/L\), so the \(j\)-th cell is
centered at \(2\pi jt/L\) (reverse the sign if using \(z=t-i\theta\)).
More directly, put \(y=\theta/t\) and \(C_n=A-L^2/2\). The exact
completed-square identity is

\[
 \Re\{A_j/(t+i\theta)+N(t+i\theta)\}
 =2s-\frac{C_n y^2}{t(1+y^2)}
 -\frac{2\pi^2}{t(1+y^2)}\left(j-\frac{Ly}{2\pi}\right)^2. \tag{12a}
\]

On the restricted interval, \(|Ly/(2\pi)|\le J+1/2\). Thus every
\(|j|\ge J+1\) has distance at least 1/2 in the final square, and
summing this Gaussian tail gives \(O(e^{2s-c_J/t})\). This justifies
replacing the infinite theta sum by the finite sum \(|j|\le J\) on
this interval. One may integrate each of these finitely many terms on
the entire restricted interval; assigning individual cells is
unnecessary. The identity also shows that the sum of absolute values of
the finite terms never has an exponentially large maximum relative to
\(e^{2s}\).

For each \(|j|\le J\) the saddle is \(z_j=\sqrt{A_j/N}\). It satisfies

\[
 \Re z_j=t\{1+O_J(L^{-4})\},\qquad
 \Im z_j=2\pi jt/L+O_J(t/L^3).
\]

Move the local vertical segment through \(\Re z_j\), then to the circle
\(|z|=|z_j|\). All its endpoints are separated from \(\arg z_j\) by a
fixed multiple of \(1/L\), so the joining contours and omitted circular
arcs lose \(c_J/t\) in real exponent. On that circle, writing
\(\beta=\arg z_j\), the phase is

\[
 A_j/z+Nz=2\sqrt{NA_j}\cos(\phi-\beta),
\]

and its real part is maximal at \(\phi=\beta\). Thus the completed
integral is precisely \(T_{j,m}\), with exponentially negligible
endpoint error. The Gaussian absolute integral is
\(O(te^{2\Re\sqrt{NA_j}}/\sqrt s)\), uniformly for \(|j|\le J\). Summing
the amplitude remainders gives \(B_0O(t^{M+1}/u)\). One can deform the
analytic remainder along with the exact small-fugacity amplitude before
truncating. Alternatively, on the original vertical line the maximum
real exponent for sector \(j\) is \(s+(\Re\sqrt{A_j})^2/t\), which is at
most 2s; its local absolute-integral width is comparable to
\(t/\sqrt s\), uniformly for the finitely many \(j\). No division by the
theta sum or exponentially growing relative-error factor is used.

A fully mechanical version of the contour argument can integrate each of
the finitely many exponentials on the entire restricted vertical
interval, then join to its saddle circle. Their endpoint losses are at
least \(c_J/t\) because every retained saddle stays a distance
comparable to \(t/L\) from the two outer endpoints. Terms whose centers
lie just outside the interval have only endpoint Gaussian tails. This
avoids making any assertion about dominance exactly at the cell
boundaries.

\subsection{Bounding the rest without an infinite Bessel
sum}\label{mfp:bounding-the-rest-without-an-infinite-bessel-sum}

Let \(C_n=A-L^2/2\). Then \(C_n\) tends uniformly to \(\pi^2/6\) and is
bounded between two positive constants. Since

\[
 A_j=\tfrac12(L+2\pi ij)^2+C_n,
\]

the real part of \(\sqrt{A_j}\) is positive and decreases with
\textbar{}\(j\)\textbar. One way to see this is to write
\(a_j=\Re\sqrt{A_j},\ b_j=\Im\sqrt{A_j}\); the equations
\(a_jb_j=\pi jL\) and \(a_j^2-b_j^2=A-2\pi^2j^2\) show that \(a_j\)
decreases from \(\sqrt A\) to \(L/\sqrt2\) when \(C_n>0\). Equivalently
differentiate the explicit formula \(a_j^2=(|A_j|+A-2\pi^2j^2)/2\). It
follows that \(\Delta_j\) increases with \textbar{}\(j\)\textbar.
Expanding the square root for fixed \(j\) gives (4).

Each of the finitely many omitted sectors \(K<|j|\le J\) is therefore
\(O(B_0e^{-\Delta_{K+1}})\). Outside the enlarged interval, (8)-(9) give

\[
 |\hbox{outer contribution}|/B_0
 \le n^{C_0}\exp\{-c_0(2J+1)^2/[tL^2]\}
\]

for all sufficiently large \(n\); the eventual condition
\((2J+1)\pi/L<1\) ensures the small-angle branch of (9) applies at the
boundary, and its monotone lower bound handles all larger angles. Choose
\(J\), fixed in \(n\), so \(c_0(2J+1)^2\) exceeds twice an eventual
upper bound for \(\Delta_{K+1}tL^2\). That upper bound exists by (4).
The polynomial factor is then absorbed, giving
\(O(B_0e^{-\Delta_{K+1}})\). All connector and modular errors
\(e^{-c_J/t}\) are much smaller than this rate. This proves (3), without
any infinite inversion or growing-\(K\) assertion.

\subsection{Exact sector functions and a genuinely sectorwise
statement}\label{mfp:exact-sector-functions-and-a-genuinely-sectorwise-statement}

To distinguish an exponential sector before truncating a much larger
principal algebraic series, define the analytic amplitude

\[
 H_{\rm ex}(z,u)=\frac{\sqrt{1+1/u}
 \exp\{-\operatorname{Li}_2(-1/u)/z+c_1(1/u)z\}}
 {(-1/u;e^{-z})_\infty}. \tag{12b}
\]

It is analytic and nonzero for \(\Re z>0,\ u>1\). Its endpoint expansion
is exactly \(\sum_mh_mz^m\). For each fixed integer \(j\) and
sufficiently large \(n\), let \(\Gamma_j\) be the vertical segment
\(z=t+i\theta\) with

\[
 (2j-1)\pi t/L\le\theta\le(2j+1)\pi t/L,
\]

oriented upward, and set

\[
 \mathcal B_j(n,u)=\frac{(-1)^j}{2\pi i\sqrt{1+u}}
 \int_{\Gamma_j} H_{\rm ex}(z,u)e^{A_j/z+Nz}\,dz. \tag{12c}
\]

These are explicit exact local integrals, not infinite formal series.
Define \(R_K\) by the exact identity
\(R_K=a_n-\sum_{|j|\le K}\mathcal B_j\). Then, for fixed \(K\),

\[
 a_n(\alpha)=\sum_{|j|\le K}\mathcal B_j(n,n^\alpha)
 +B_0 O_{a,b,K}(e^{-\Delta_{K+1}}). \tag{12d}
\]

For each fixed \(j\) and \(M\),

\[
 \mathcal B_j(n,u)=\frac{(-1)^j}{\sqrt{1+u}}
 \sum_{m=0}^{M}h_m(u)T_{j,m}
 +B_0e^{-\Delta_j}O_{a,b,j,M}(t^{M+1}/u). \tag{12e}
\]

For \(j=0\) use \(\Delta_0=0\). Proof of (12d): use the auxiliary \(J\)
from Section 5.3. Inside its full arc, (11) is exactly the sum of the
\(H_{\rm ex}\) sector integrands times \((Q;Q)_\infty^{-1}\). The
modular factor differs from 1 by \(e^{-c/t}\). For the \(j\)-th
exponential, the contribution outside its own \(\Gamma_j\) but inside
the full arc is \(e^{-c_J/t}\) times the principal scale, by (12a),
because \(|j-Ly/(2\pi)|\ge1/2\). All infinitely many exterior \(\theta\)
indices have the same summable Gaussian bound. This leaves precisely the
finite sum of exact local integrals \(|j|\le J\) plus \(e^{-c_J/t}\)
errors. Complete each fixed local integral to its saddle contour to
bound its modulus by \(O(B_0e^{-\Delta_j})\); then drop \(K<|j|\le J\)
and use the already established outer-arc estimate. This proves (12d).

For (12e), first deform the exact \(H_{\rm ex}\) integrand to the short
saddle-circle arc through \(z_j\), keeping the entire deformation in
\(\Re z\) comparable to \(t\) and \(|\Im z|\le C_jt/L\). Apply and
truncate the uniform small-fugacity Euler--Maclaurin expansion there.
Only the resulting polynomial integrands are then completed around the
full circle and evaluated by (10); \(H_{\rm ex}\) is never continued
into \(\Re z\le0\). On the short saddle arc the Gaussian
absolute-integral scale is \(B_0e^{-\Delta_j}\), so the analytic
amplitude remainder has the required sectorwise bound. Endpoint and
connector errors lose \(c_j/t\) in real phase and are smaller than
\(t^{M+1}/u\) relative to the sector scale. This proves (12e). The order
of these steps matters: on the original vertical line its maximum real
phase exceeds \(2\Re\sqrt{NA_j}\) by \((\sqrt A-\Re\sqrt{A_j})^2/t\),
which need not be bounded. Estimating the sectorwise error there would
be insufficient.

Thus the exponential error in (12d) is genuine before truncating the
principal sector, and (12e) gives the algebraic expansion separately
within every fixed sector. Numerically evaluating the exact local
integrals is an alternative to using finite-\(M\) Bessel models; no
implementation of certified integral quadrature is claimed in the
accompanying code.

\Needspace{9\baselineskip}
\section{OEIS specializations and the fugacity
threshold}\label{mfp:oeis-specializations-and-the-fugacity-threshold}

The dilogarithm identity gives

\[
 A=A_\infty+\operatorname{Li}_2(-1/u),\qquad
 A_\infty=\tfrac12\alpha^2\log^2n+\pi^2/6,
\]

and \(\operatorname{Li}_2(-1/u)=-1/u+O(u^{-2})\). Therefore the fully
elementary Sommerfeld expression

\[
 E_n(\alpha)=\frac{A_\infty^{1/4}}{2\sqrt\pi\,n^{\alpha/2+3/4}}
 e^{2\sqrt{nA_\infty}}
\]

satisfies \(a_n/E_n\to1\) \textbf{if and only if \(\alpha\ge1/2\)}, for
fixed \(\alpha>0\). At \(\alpha=1/2\) the omitted exponent still tends
to zero, but only at rate \(1/\log n\). For \(0<\alpha<1/2\),

\[
 \log\frac{a_n}{E_n}
 \sim-\frac{\sqrt n}{u\sqrt{A_\infty}}
 \sim-\frac{\sqrt2}{\alpha}\frac{n^{1/2-\alpha}}{\log n},
\]

so \(a_n/E_n\) tends to zero. This is a relative-asymptotic issue, not a
small change to the leading logarithmic scale.

For \(\alpha=1\), the theorem proves the conjectural equivalent
displayed for A291698:

\[
 a_n\sim\frac{(\log^2n+\pi^2/3)^{1/4}}
 {\sqrt\pi(2n)^{5/4}}
 \exp\sqrt{2n(\log^2n+\pi^2/3)}.
\]

For \(\alpha=2\) it proves the displayed A292304 equivalent:

\[
 a_n\sim\frac{(2\log^2n+\pi^2/6)^{1/4}}{2\sqrt\pi n^{7/4}}
 \exp\{2\sqrt{n(2\log^2n+\pi^2/6)}\}.
\]

Keeping A exact, N shifted, and the factor \(\sqrt{1+u}\) exact is the
clean route to the uniform theorem and its inverse.

\section{Real continuation and
inversion}\label{mfp:real-continuation-and-inversion}

A sequence has no unique real analytic continuation. To make the inverse
unambiguous, specify the following \textbf{saddle-Fourier continuation},
for real \(x\) sufficiently large:

\[
 \mathcal A_\alpha(x)=\frac1{2\pi}
 \int_{-\pi}^{\pi}e^{x(t(x)-i\theta)}
 F(e^{-t(x)+i\theta},x^\alpha)\,d\theta, \tag{13}
\]

where A,N,\(t\) are defined as above with \(n\) replaced by \(x\). This
integral is real, and equals \(a_n(\alpha)\) at every positive integer
\(x=n\). This choice is a convention; no claim of uniqueness among
continuations is made.

The preceding major/minor estimates work for real \(x\) without
alteration, since \(|e^{-ix\theta}|=1\). In particular (13) is positive
eventually. They can also be differentiated a fixed number of times:
differentiating products introduces fixed sums of k powers and factors
\(u'/u,\ t',\text{ and }\theta\), all bounded by polynomial factors in
\(x\) and log(\(x\)); removing a fixed number of product factors changes
the minor-arc damping by at most a fixed constant. In the main arc
differentiate the small-fugacity Euler--Maclaurin representation, taking
additional endpoint depth as needed. More explicitly, let
\(R_P(x,\theta)\) be the small-fugacity Euler--Maclaurin remainder used
in Section 5.1, with \(\theta\) held fixed. For every fixed \(d\),

\[
 \partial_x^d R_P(x,\theta)
 =O_{a,b,P,d,C}\left(\frac{t^{2P-1}}{u x^d}\right),
 \qquad |\theta|\le Ct/L. \tag{13a}
\]

This follows directly by differentiating its periodic-Bernoulli
remainder integral: \(v=1/u\) has \(v^{(r)}=O(vx^{-r})\),
\(t^{(r)}=O(tx^{-r})\), and the occurring weighted derivative integrals
satisfy

\[
 \int_0^\infty y^k
 \left|\partial_w^\ell D_v^h\log(1+v e^{-w})\big|_{w=yz}\right|dy
 =O(v/t^{k+1})
\]

for the nontrivial derivatives that occur. Factors from differentiating
\(z^{2P}\) and from \(yt'(x)\) supply precisely the displayed powers.
Moving arc endpoints contribute only \(e^{-c/t}\)-scale boundary terms.
On minor arcs, product differentiation removes at most \(d\) factors and
costs only fixed polynomial factors. Consequently errors can still be
made smaller than any prescribed algebraic order after any fixed number
of derivatives; the unchanged exact \(\Delta_{K+1}\) rate is not
asserted for these differentiated errors. In particular,

\[
 G_\alpha(x)=\log\mathcal A_\alpha(x),\qquad
 G_\alpha'(x)\sim\frac{\alpha(\log x+2)}{\sqrt{2x}}>0. \tag{14}
\]

For an explicit check of the first derivative, write \(z=t-i\theta\),
\(F=F(e^{-z},u)\), and differentiate at fixed \(t\) before accounting
for \(t=t(x)\). Integration by parts in \(\theta\) gives the exact
identity

\[
 \mathcal A_\alpha'(x)=\frac1{2\pi}\int_{-\pi}^{\pi}e^{xz}
 \left\{zF+\frac{\alpha u}{x}F_u\right\}\,d\theta
 +\frac{t'(x)}\pi e^{xt}F(-e^{-t},u)\sin(\pi x). \tag{14a}
\]

The last term is a minor-arc boundary term and is beyond every power
relative to the main term. On the principal arc, the leading multiplier
is \(z+\alpha\log(1+u)/(xz)\), since \(uA'(u)=\log(1+u)\). Its saddle
value is \(t+\alpha\log(1+u)/(xt)\), giving (14). This argument uses the
derivative of the product itself on the minor arcs, so possible zeros of
F cause no division problem.

Thus (13) is eventually strictly increasing and has a well-defined large
positive inverse.

If X is the target and \(y=\log X\), set

\[
 z=\frac{y}{2\sqrt2\alpha},\qquad r_0=\frac z{W_0(z)},\qquad x_0=r_0^2. \tag{15}
\]

Then \(\mathcal A_\alpha^{-1}(X)\) is asymptotic to \(x_0\). Writing
\(l_0=\log x_0\), the first explicit inverse correction is

\[
 \mathcal A_\alpha^{-1}(X)
 =x_0-\frac{\pi^2x_0}{3\alpha^2l_0(l_0+2)}
 +O_{a,b}(x_0/l_0^4). \tag{15a}
\]

Indeed
\(G(x)-g(x)=\frac{\sqrt2\pi^2}{6\alpha}\frac{\sqrt x}{\log x}+O(\sqrt x/\log^3x)+O(\log x)+O(x^{1/2-\alpha}/\log x)\).
Divide the first term by g'(\(x\)), and apply Taylor's theorem. The
quadratic displacement and the next phase term are \(O(x/\log^4x)\); the
prefactor and finite-\(u\) corrections are smaller uniformly for
\(\alpha\ge a>0\). This is a first logarithmic inverse transseries term,
not an absolute-error certificate.

For controlled finer corrections let \(S_{K,M}(x)\) be the finite sum in
(3), and solve

\[
 \log S_{K,M}(x)=y
\]

by Newton iteration beginning near \(x_0\). The approximant is positive
and increasing eventually; its derivative has the same leading
equivalent (14). If \(x_*\) is its large solution and \(x\) is the exact
inverse of (13), then

\[
 x-x_*=O\left(\frac{\sqrt{x_*}}{\log x_*}
 \left\{\frac{t(x_*)^{M+1}}{x_*^\alpha}+e^{-\Delta_{K+1}(x_*)}\right\}\right). \tag{16}
\]

This is an arbitrary-fixed-algebraic-order inverse procedure when \(M\)
is increased. For \(\alpha\) below 1/2, the exact dilogarithm must
remain in S; replacing it by \(A_\infty\) before inversion discards an
unbounded logarithmic correction.

A formal inverse coefficient generator is also available. Let
\(g(x)=\sqrt2\alpha\sqrt x\log x\), let \(H=g^{-1}\) be given by (15),
and put \(E(x)=\log S_{K,M}(x)-g(x)\), \(\eta(y)=E(H(y))\). Lagrange
reversion in an auxiliary parameter multiplying E gives

\[
 x=H(y)+\sum_{r\ge1}\frac{(-1)^r}{r!}
 \left(\frac d{dy}\right)^{r-1}\{H'(y)\eta(y)^r\}. \tag{17}
\]

Equation (17) is a formal reversion generator; truncating it requires
the usual ordered error bookkeeping. A numerical Newton solve with (16)
avoids asserting convergence of an infinite transseries. Every
derivative includes \(u=x^\alpha\), \(A(u)\), \(N(x)\), and the
finite-sector dependence.

\subsection{Integer threshold}\label{mfp:integer-threshold}

The distinct-partition map that increments the largest part embeds
partitions of \(n\) into those of \(n\)+1 without changing part count.
Since \((n+1)^\alpha>n^\alpha\), the weighted sequence is strictly
increasing for \(n\ge1\). For sufficiently large X, because (13) is
monotone and agrees with the sequence at integers,

\[
 \min\{n:a_n(\alpha)\ge X\}=\left\lceil\mathcal A_\alpha^{-1}(X)\right\rceil.
\]

An approximate inverse must \textbf{not} simply be rounded when its
error interval crosses an integer. Use an enclosure from a certified
remainder, or evaluate the neighboring exact coefficients. The big-O
constants in this note are not numerical certificates for a particular
finite input.

\begin{writenote}[An instance of the repository's inversion apparatus]
No novelty is claimed here for the inversion method; the instances are
those of the transseries volume
(\nolinkurl{Analysis/Transseries/docs/series-and-transseries/Transseries_And_Inversion/transseries_and_inversion.tex}).
The initializer (15) solves $r\log r=z$, that is $X+\log X=\log z$ for
$X=\log r=\tfrac12\log x$: the dominant block of
\texttt{p0:thm:lambert-core} with $a=b=1$, on the principal branch, so
$r_0=e^{W_0(z)}=z/W_0(z)$. The generator (17) is the reversion formula of
\texttt{p0:thm:perturbed-inversion}, read formally at $\varepsilon=1$ with
$F=g$, $\mu_0=H$ and $E(\mu_0(y))=\eta(y)$; the manuscript itself calls it
formal and does not assert convergence. The threshold identity of this
subsection is part~(1) (exact rounding) of \texttt{p0:thm:staircase} for the
interpolation (13), and the warning against rounding an approximate inverse
is its separation condition, part~(2).
\end{writenote}

\section{Exact computations and practical
warning}\label{mfp:exact-computations-and-practical-warning}

The code \texttt{checks.py} computes the exact polynomials

\[
 [q^n]F(q,u)=\sum_{m\ge1}p_{\le m}(n-m(m+1)/2)u^m
\]

using the standard restricted-partition dynamic program. For
\(\alpha=1,2\) the evaluation is exact integer arithmetic, followed by
100-digit comparisons. For fractional \(\alpha\) the polynomial
coefficients remain exact and evaluation uses 100-digit arithmetic. The
displayed low \(\alpha=2\) values agree term for term with OEIS through
\(n=23\); \(\alpha=1\) gives 1,1,2,12,20,55,294,497,1224,\ldots{} .

At \(n\)=20,000:

\begin{itemize}
\tightlist
\item
  \(\alpha=1\): principal Bessel relative error is about \(-9.0139\times10^{-5}\)
\item
  \(\alpha=2\): principal Bessel relative error is about \(-5.5771\times10^{-2}\)
\item
  \(\alpha=1/4\): principal relative error is about \(-2.9187\times10^{-10}\), reduced
  to about \(2.24\times10^{-16}\) by the \(c_2\) correction
\end{itemize}

Adding \(c_2,c_3,\ldots\) does essentially nothing to the first two
errors because the oscillatory sectors dominate the algebraic endpoint
correction at these sizes. This late onset is explained, rather than
contradicted, by (4).

\texttt{resonance\_check.py} compares finite conjugate Bessel sectors,
keeping endpoint degree \(M=9\). At \(n\)=20,000:

\begin{itemize}
\tightlist
\item
  \(\alpha=1\): \(K=0\) gives \(-9.0139\times10^{-5}\), \(K=1\) gives \(+4.43\times10^{-9}\),
  \(K=2\) gives \(-4.92\times10^{-12}\)
\item
  \(\alpha=2\): \(K=0\) gives \(-5.5771\times10^{-2}\), \(K=1\) gives \(-1.0985\times10^{-2}\),
  \(K=2\) gives \(-5.09\times10^{-5}\)
\end{itemize}

These numerical improvements are diagnostics, not certified finite-\(n\)
error bounds. Increasing \(K\) indefinitely is unsafe: (3) fixes \(K\)
before \(n\) tends to infinity, and the formal infinite Bessel sum need
not converge. Remote root-of-unity arcs and modular factors in (11) can
matter before the eventual regime. The script records all \(K\) from 0
through 8 so this limitation is visible rather than hidden.

\section{Prior-art boundary and sources
checked}\label{mfp:prior-art-boundary-and-sources-checked}

\begin{itemize}
\tightlist
\item
  OEIS A291698: \url{https://oeis.org/A291698}. Sequence definition and leading conjecture.
\item
  OEIS A292304: \url{https://oeis.org/A292304}. Sequence definition and leading conjecture.
\item
  Vivien Brunel, \emph{A general asymptotic formula for distinct
  partitions}, arXiv:1709.04955, journal DOI
  10.1016/j.aop.2018.03.023: \url{https://arxiv.org/pdf/1709.04955}
\item
  Dan Romik, \emph{Partitions of \(n\) into \(t\sqrt n\) parts},
  European Journal of Combinatorics 26 (2005), 1--17, DOI
  10.1016/j.ejc.2004.02.005:
  \url{https://kam.mff.cuni.cz/~ksemweb/clanky/Romik-partitions.pdf}
\item
  Arash Arabi Ardehali and Hjalmar Rosengren, \emph{A New Product
  Formula for \((z;q)_\infty\), with Applications to Asymptotics},
  published 19 September 2026:
  \url{https://link.springer.com/article/10.1007/s00365-026-09783-2}
\item
  Richard J. McIntosh, \emph{Some asymptotic formulae for q-shifted
  factorials}, Ramanujan Journal 3 (1999), 205--214:
  \url{https://doi.org/10.1023/A:1006949508631} (bibliography verified in the
  2026 primary paper; this 1999 full text was not reviewed here)
\item
  NIST DLMF: Jacobi triple product (Chapter 17), \(\theta\)
  transformations (Chapter 20), and modified Bessel large-argument
  asymptotics (10.40): \url{https://dlmf.nist.gov/17.8} ;
  \url{https://dlmf.nist.gov/20.7} ; \url{https://dlmf.nist.gov/10.40}
\end{itemize}

Brunel's Section 3 treats prescribed part count by a two-variable
saddle and recovers classical regimes; it does not state the polynomial
moving-fugacity all-order remainder or finite conjugate-sector theorem
above. Romik's Theorem 2 is uniform for the normalized part count on
compact subsets strictly below the distinct-partition endpoint
\(\sqrt2\); polynomial fugacity approaches that endpoint. His
local-limit proof is relevant methodological prior art. These checks
delimit those two sources; they do not prove a global novelty claim, and
the grand-canonical/Fermi-gas interpretation and leading saddle method
are classical.

The September 2026 paper supplies a uniform moving-argument q-product
expansion (Corollary 2) and discusses the classical large-argument
theta regime. It does not state the coefficient-diagonal or inverse
theorem studied here. Its Corollary 2 assumes a fixed argument for the
small complex parameter and a lower bound on the real part of the
transformed argument. These are relevant antecedents, not a basis for
claiming that q-product Euler--Maclaurin or modular transformations are
new.

\begin{writenote}[The two OEIS conjectures]
The delivered proof note (\nolinkurl{theorem-and-proof.md}, Section~9)
records that the OEIS page of A291698 could not be reached when the
manuscript was written, so that its displayed conjecture was transcribed
from a screening note that is not part of the delivery, while the A292304
entry was checked directly. The manuscript's list above drops that
caveat. On 2 October 2026 both entries were consulted for this write: the
equivalents displayed in Section~6 are exactly the conjectures of V\'aclav Kot\v{e}\v{s}ovec in A291698
(15 September 2017) and A292304 (14 September 2017), which the manuscript
does not name; and the listed terms $1,1,2,12,20,55,294,497,1224,2520$ of
A291698 agree with an independent recomputation.
\end{writenote}

\begin{writenote}[Relation to other reports of the collection]
The neighbouring report \texttt{a022629-\allowbreak distinct-\allowbreak
partition-\allowbreak norms}, in the same directory as this one, treats
the norm-weighted products $\prod_{k\ge1}(1+k^\alpha q^k)^\mu$ (A022629 and
its relatives), where the weight $k^\alpha$ depends on the \emph{part}
$k$. Here the weight $n^\alpha$ depends on the \emph{size} $n$ being
counted and is the same for every part, so a partition with $m$ parts
carries $n^{\alpha m}$. The two families share the distinct-part structure,
the Euler--Maclaurin and saddle methods and a Lambert-type inverse, but
not the saddle regime (here the part count approaches the
distinct-partition endpoint $\sqrt{2n}$, with conjugate resonance sectors
governed by $L=\log u$), the sequences or any theorem. Neither manuscript cites the other. That report
was not edited by this write.
\end{writenote}

\section{Further research questions}
\subsection{Effective constants and certified numerical ranges}
The present bounds are uniform but eventual. An explicit value for the consecutive-interval constant, quantitative Euler--Maclaurin remainders, and certified saddle-contour quadrature would turn the inverse enclosure into a finite-input certificate. This is especially relevant for $\alpha=2$, where asymptotic smoothness begins late. The exact local integrals in (12c) provide a natural numerical starting point; the supplied quadrature experiment is not interval arithmetic.

\subsection{A growing number of resonance sectors}
The proof fixes $K$ and an auxiliary $J$ before sending $n$ to infinity. Extending it to $K=K(n)$ requires uniform control of the changing saddle geometry and the modular factor $(Q;Q)_\infty^{-1}$. Near remote root-of-unity arcs that factor need not be close to one. An unqualified infinite sum of the formal Bessel sectors is not justified and may fail to converge.

\subsection{The transition beyond polynomial fugacity}
The first resonance scale is governed by $\sqrt n/L^3$, where $L=\log u$. Polynomial fugacity has $L\asymp\log n$, so every fixed nonprincipal sector is beyond all algebraic orders. A different regime should appear when $L$ is comparable to $n^{1/6}$, for which that damping is of order one. The present theorem does not cover that regime. A uniform theta-type crossover would need its own proof rather than an assumption that a Gaussian approximation remains accurate for a small part-count variance.

\subsection{Comparison with general q-product asymptotics}
The endpoint expansion and modular ingredients are classical. The appropriate contribution to compare against the literature is the diagonal coefficient transfer with moving fugacity, explicit finite exponential sectors, and controlled inversion. The recent moving-argument product formulas cited in Section 9 may simplify effective remainder calculations. A broader search is still needed before making a publication-level novelty claim.

\clearpage
\appendix
\section{Reproducibility and independent checks}
The reproducibility package contains the complete editable \LaTeX\ source, the proof-note source, four Python programs, and their saved numerical output. The mathematical source used for integration has SHA256
\begin{quote}\small\ttfamily
5a90f1336fbbd5c0804f4781f436d47cb18a47ef82cefc04ef1f322d5cc9f64b
\end{quote}
The long hash identifies the source precisely; it is not a mathematical certificate.

\subsection{Programs and outputs}
\begin{description}[style=nextline,leftmargin=1em,labelindent=0pt,labelwidth=0pt]
\item[\texttt{checks.py}] Generates exact restricted-partition polynomials, the endpoint coefficient recurrence, principal Bessel comparisons through $n=20{,}000$, and principal-model inverse checks. \\\emph{Output: }\texttt{checks.json}.
\item[\texttt{resonance\_check.py}] Evaluates the finite conjugate Bessel sums for $K=0,\ldots,8$, with $M=9$. \\\emph{Output: }\texttt{resonance-checks.json}.
\item[\texttt{validate.py}] Performs 120 independent direct weighted-product dynamic-programming checks for $n\le60$, verifies 24 listed A292304 terms, tests the exact Jacobi identity at real and complex arguments, and compares finite-sector inverses. \\\emph{Output: }\texttt{validation.json}.
\item[\texttt{sector\_integrals.py}] Performs uncertified high-precision quadrature of the exact local integrals at $\alpha=1$, $n=20{,}000$, for $j=0,1,2,3$. \\\emph{Output: }\texttt{sector-integral-checks.json}.
\end{description}

\begin{writenote}[Shipped names]
In this repository the four programs are \nolinkurl{code/checks.py},
\nolinkurl{code/resonance_check.py}, \nolinkurl{code/validate.py} and
\nolinkurl{code/sector_integrals.py}, their outputs are in
\nolinkurl{data/}, the delivered \nolinkurl{build_pdf.sh} is
\nolinkurl{code/build_pdf.sh}, and \nolinkurl{moving-fugacity-report.tex} is
\nolinkurl{article.tex}. The last three programs read and write files
beside themselves, so the replay below must be run on a flat copy; the
report's README gives the commands.
\end{writenote}

The checks were run with Python, mpmath 1.3.0, and SymPy 1.14.0. The principal comparison uses 100 decimal digits; the local-sector quadrature uses 70. All integer polynomial coefficients are exact. No random seed is involved.

\subsection{Clean replay}
From the package directory, run
\begin{verbatim}
python checks.py --max-n 20000 --output checks.json
python resonance_check.py
python validate.py
python sector_integrals.py
pdflatex -interaction=nonstopmode moving-fugacity-report.tex
pdflatex -interaction=nonstopmode moving-fugacity-report.tex
\end{verbatim}
The JSON values are high-precision numerical observations, not rigorously rounded enclosures. The proof does not depend on their being small.

\Needspace{12\baselineskip}
\subsection{Finite-sector inverse diagnostics}
For target $y=\log a_{20000}(\alpha)$, Newton inversion of the $M=9$ model gives the following signed errors $x_*-20000$:
\begin{center}
\begin{tabular}{rrrr}
\toprule
$\alpha$ & $K=0$ & $K=1$ & $K=2$\\
\midrule
1 & $1.49946\times10^{-3}$ & $-7.36613\times10^{-8}$ & $8.18152\times10^{-11}$\\
2 & $4.81105\times10^{-1}$ & $8.31548\times10^{-2}$ & $3.80911\times10^{-4}$\\
\bottomrule
\end{tabular}
\end{center}
These values demonstrate the importance of the first conjugate sectors. They do not certify a ceiling when the true inverse lies at an integer: the exact coefficient or an error enclosure must still decide the threshold.

\subsection{Identity and local-integral diagnostics}
Three direct evaluations of the exact Jacobi identity had relative discrepancies below $3\times10^{-77}$ at the selected working precision. Independent direct weighted-product evaluation agreed exactly with the restricted-partition-polynomial method in all 120 small cases. All 24 listed A292304 terms agreed exactly.

At $\alpha=1$, $n=20{,}000$, the $M=9$ expansion matched the numerically integrated principal local sector to about $1.28\times10^{-28}$ relative, and the first complex sector to about $3.61\times10^{-27}$. Higher fixed cells at this finite size had larger endpoint effects; no monotonic improvement in sector count is promised. These checks support the coefficient signs and normalizations but are not substitutes for the contour proof or certified quadrature.

\section{Scope checklist}
\begin{itemize}
\item Uniformity is on a fixed positive compact range of $\alpha$.
\item The sector count $K$, auxiliary count $J$, and endpoint depth $M$ are fixed before $n$ tends to infinity.
\item The exact amplitude is deformed only within $\Re z>0$; only truncated polynomial integrands are completed around a full circle.
\item Exact dilogarithms are retained for general positive $\alpha$; the elementary Sommerfeld equivalent is relatively valid only for fixed $\alpha\ge1/2$.
\item The real inverse refers to the explicitly defined saddle--Fourier continuation.
\item Big-$O$ constants in the report are not finite-input numerical certificates.
\item No external submission, OEIS edit, or publication-priority assertion is part of this report.
\end{itemize}
\end{document}
